\chapter{Jumps, representatives, and boundary traces}\label{sec:examples-bv-interval}
These examples supplement the singular-continuous model in Example \ref{ex:cantor-function-bv}. They use the exact interval representation of Section \ref{sec:bv-interval-structure} to distinguish a function's almost-everywhere class, its canonical representative, and its endpoint traces.

\begin{example}[A BV function with dense jumps]\label{ex:bv-dense-jumps}
Enumerate the rationals in $(0,1)$ without repetition as $q_1,q_2,\ldots$ and set
\[
 u(x)=\sum_{j=1}^\infty 2^{-j}\chi_{[q_j,1)}(x),\qquad 0<x<1.
\]
Then $u$ is right-continuous, strictly increasing and in $BV(0,1)$, with
\[
 Du=\sum_{j=1}^\infty 2^{-j}\delta_{q_j},\qquad |Du|((0,1))=1.
\]
It jumps by $2^{-j}$ at $q_j$ and is continuous at every irrational point. Its endpoint traces are zero and one.
\end{example}
\begin{proof}
The series converges uniformly by the summable bounds. The measure $\mu=\sum_j2^{-j}\delta_{q_j}$ is a Radon probability measure by the summable-tail argument of Example \ref{ex:dense-atomic-measure}, now enumerating only interior rationals. The displayed function is its cumulative function; Lemma \ref{lem:interval-measure-primitive} gives its derivative, one-sided continuity and traces. Every nonempty interval contains a rational carrying positive mass, so every increment between distinct points is strictly positive. The atom formula in Theorem \ref{thm:interval-bv-representative} gives the jumps. At an irrational point the measure has no atom, so the left and right limits agree. Thus dense discontinuities are compatible with finite total variation; their amplitudes must be summable.
\end{proof}

\begin{example}[Pointwise variation depends on the representative]\label{ex:dirichlet-function-bv-representative}
The function $u=\chi_{\mathbb Q\cap(0,1)}$ is zero almost everywhere, belongs to $BV(0,1)$ and has distributional variation zero. As a specified pointwise function it has infinite classical variation. Its canonical right-continuous BV representative is identically zero.
\end{example}
\begin{proof}
Countability makes the rationals Lebesgue null, so all distributional pairings agree with those of zero. To obtain arbitrarily large pointwise variation, choose an arbitrarily long increasing sequence alternating between rational and irrational points, possible by density of both sets. Every successive function increment has absolute value one. The regulated representative theorem applies to the almost-everywhere class and gives the zero representative; it does not assign finite classical variation to every null modification of that class.
\end{proof}

\begin{example}[A boundary layer can retain a trace while disappearing in the interior]\label{ex:bv-boundary-layer-trace}
Let $u_j=\chi_{(0,1/j)}$ on $I=(0,1)$, $j\geq2$. Then $u_j\to0$ in $L^1(I)$, $|Du_j|(I)=1$, and $Du_j\rightharpoonup0$ against compactly supported continuous tests in $I$. Yet $u_j(0+)=1$ for every $j$. The zero extensions to $\mathbb R$ have derivative $\delta_0-\delta_{1/j}$ and total variation two.
\end{example}
\begin{proof}
The $L^1$ norm is $1/j$. Within the open interval there is only the right endpoint, so $Du_j=-\delta_{1/j}$ by Theorem \ref{thm:interval-finite-perimeter-sets}. Every compact support stays away from zero and eventually misses that atom. The interior right trace at zero is one. Proposition \ref{prop:interval-bv-zero-extension} adds the missing boundary atom $\delta_0$, whose support is distinct from $1/j$. The example shows that an $L^1$ limit with uniformly bounded variation does not alone preserve boundary traces.
\end{proof}

\begin{example}[Three derivative types in one BV function]\label{ex:bv-three-derivative-types}
Let $C$ be the Cantor distribution function on $[0,1]$, and let $0<c<1$. The function $u(x)=x+C(x)+\chi_{[c,1)}(x)$ on $(0,1)$ is in $BV(0,1)$ and
\[
 Du=\mathcal L^1\lfloor(0,1)+\mu_C\lfloor(0,1)+\delta_c,
 \qquad |Du|((0,1))=3.
\]
Its absolutely continuous, singular-continuous and atomic summands are mutually singular, even when $c$ belongs to the Cantor set.
\end{example}
\begin{proof}
Differentiate each summand in distributions, using Example \ref{ex:cantor-function-bv} and the interval endpoint formula. The Cantor measure is singular to Lebesgue measure and has no atoms, including at the endpoints and at $c$. All three measures are positive of mass one, so their sum has total variation three. Nonatomicity lets one remove $c$ from a full-measure support of the Cantor part to obtain disjoint carriers for the three types. This also shows why absence of jumps does not by itself imply absolute continuity.
\end{proof}
